\scanpage{117}
% archive sheet number: 116
point est obtenu, pour un $i$ et un $z_i$ fixé, par le choix d'un $z'_i$ satisfaisant
\[
 {z'_i}^{\,2}=\underbrace{z_i^3+a_1z_i^2+a_2z_i+a_3}_{(z_i-\unsure{\zeta_i|\xi_i})(z_i-\zeta_j)(z_i-\zeta_k)}. \tag*{(66)}
\]
Choisir, pour tt $i\in\unsure{I|J}$, un $x_i\in E(k)$ tel que
\[
 2x_i=\varepsilon_i \tag*{(67)}
\]
revient donc à choisir, pour tt $i$, \deleted{\unsure{des|les}} la quantité
\[
 \wp(x_i)=z_i \tag*{(68 a)}
\]
soumise à (65) \inserted{(2 choix possibles)}, et ensuite, la quantité $z'_i$
\[
 \wp'(x_i)=z'_i) \tag*{(68 b)}
\]
soumise à (66) -- ce qui fait quatre choix possibles pour tt $i$, comme il se doit (car l'indétermination pour $x_i$ est dans $\vphantom{E}_2E(k)$).

Ceci posé, la condition d'épinglage de Jacobi relatif d'échelon $4$
\[
 \sum_{i\in I}x_i=0 \tag*{(69)}
\]
s'exprime par la condition
\[
 \left\{\begin{aligned}
 &\det\begin{pmatrix}1&z_1&z'_1\\1&z_2&z'_2\\1&z_3&z'_3\end{pmatrix}=0\qquad\text{i.e.}\\
 &(z_1z'_2-z_2z'_1)+(z_2z'_3-z_3z'_2)+(z_3z'_1-z_1z'_3)=0.
 \end{aligned}\right. \tag*{(70)}
\]
(cf prop. \unsure{33|3)}, qui s'applique, car les $x_i$ sont tous distincts. --)
